\section{Discussion}\label{sec:discussion}

\paragraph{RQ1: governance as compilation.}
The construct inventory shows that identity, residency, provider, policy,
governance, threat, and evidence relationships can be represented in one
language and rejected before execution when internally inconsistent.  The
benefit is not that declarations become true; it is that assumptions become
typed, referentially checked, and available to both runtime and evidence.

\paragraph{RQ2: constrained request planning.}
The compiler/VM pipeline and provider split place validation, policy,
capability, and registry checks before planning.  The observed traces contain
blocked external attempts and simulated-only execution.  The result supports a
fail-closed application architecture, while leaving host isolation untested.

\paragraph{RQ3: evidence and offline consistency.}
Trace, SecurityReport, and EvidenceBundle serve different functions: event
history, interpreted runtime state, and content linkage among bytecode, trace,
security report, and the ledger digest derived from trace events.  The report's
LedgerSummary stores that digest, not the events.  All complete bundles verify offline;
the upstream source is not part of that verification.  The policy contradiction proves why these artifacts should
not be collapsed: valid evidence can faithfully preserve an unfavorable
result.

\paragraph{RQ4: demonstrated and undemonstrated behavior.}
The snapshot demonstrates repeatable artifact production for 27 complete
sessions, six visible incomplete sessions, stable control events, and offline
digest verification.  It does not demonstrate prompt robustness, diverse
workloads, live providers, identity authentication, credential verification,
cryptographic handshakes, production containment, or compliance.

\paragraph{RQ5: open agentic-web placement.}
Project NANDA provides a research direction for indexed, verifiable agent facts
\citep{projectnanda2026,raskar2025beyonddns}; DCP-AI frames digital citizenship
and governance \citep{naranjo_dcpai}; ATrust supplies an unpublished conceptual
trust vocabulary \citep{vazquez_atrust}.  \Argorix can compile local metadata
that might feed such systems, but it is not their resolver, registry, or trust
authority.

This positioning differs from protocol conformance.  MCP and A2A define
interaction protocols \citep{mcp_spec_2025,a2a_spec}; \Argorix presently checks
declarative bridge boundaries with execution disabled.  Similarly, DID and VC
standards specify interoperable representations
\citep{w3c_did_core_2022,w3c_vc_data_model_2025}; the language stores
descriptions without resolving or verifying them.

The policy contradiction is also a design lesson for evidence-consuming
systems.  ``Passed'' should not be a free string replicated at multiple
levels.  A typed lattice such as \emph{deny} $>$ \emph{review} $>$
\emph{allow}, with unknown treated as review or deny, would make aggregation
monotone.  An EvidenceBundle should then bind both component results and the
derived decision, enabling a verifier to recompute the aggregate.

Finally, the repeated 265-event trace suggests that exhaustive metadata can
become noisy.  Evidence maps and report summaries can improve navigation, but
summary must remain traceable to component events.  Compact evidence is useful
only if it cannot erase negative findings.
